%!TEX root = main.tex

formalization of firrtl transformation

\subsection{``Expand Whens''}

HiFIRRTL allows to define conditional components,
mostly connect statements that depend on a condition.
However, a circuit cannot contain conditional connections;
they need to be replaced by multiplexers.
The ``expand whens'' transformation does exactly this:
it replaces connects within \textit{when} statements by multiplexers where necessary.
For example, \pagebreak[1]\\
\begin{minipage}[t]{12.97em}
\vspace*{-\baselineskip}
\begin{lstlisting}
module TestCondition
  input in : UInt<32>
  input condition : UInt<1>
  output pos : UInt<32>
  output neg : SInt<33>
  pos <= 0
  when condition :
    pos <= in
  else :
    wire a : SInt<33>
    a <= neg in
    neg <= a
\end{lstlisting}
\end{minipage}
is translated to\qquad{}%
\begin{minipage}[t]{16.33em}
\vspace*{-\baselineskip}
\begin{lstlisting}
module TestCondition
  input in : UInt<32>
  input condition : UInt<1>
  output pos : UInt<32>
  output neg : SInt<33>
  wire a : SInt<33>
  pos <= mux(condition, in, 0)
  a <= neg in
  neg <= validif(not condition, a)
\end{lstlisting}
\end{minipage}

Additionally, ``expand when'' removes repeated connects,
thereby implementing the \emph{last connect} semantics.
The ``expand when'' transformation produces a map
where every connected component is mapped to the one expression
that it is connected to.
This map can then easily changed back into a list of connect statements.

The specification of the Expand Whens transformation is therefore:
\begin{description}
\item[The input] of Expand Whens is the code of a module
where all types of components have been lowered to ground types.
For example, a statement \lstinline!reg a : { a1 : UInt<32>, a2 : SInt<16> },clk!
is already translated to the two statements
\lstinline!reg a$a1 : UInt<32>, clk! and
\lstinline!reg a$a2 : SInt<16>, clk!.
Also, all partial connects (\lstinline!<-!) have been replaced by appropriate full connects (\lstinline!<=!).

\item[The output] of Expand Whens is a pair (list of statements, map of connects);
the list of statements does not contain any connect statements
(but mainly declarations of registers and wires and definitions of nodes,
copied verbatim from the input),
while the map of connects contains
the last connect found on each path through the input statement sequence that declares the component;
if multiple execution paths declare the component but provide different definitions,
then a suitable multiplexer or \lstinline!validif! expression is returned.
\end{description}

In the above example, the resulting code always connects \lstinline!a <= in!
because wire \lstinline!a! is only declared in the \lstinline!else! path.
On the other hand, output \lstinline!neg! is connected using a \lstinline!validif! statement
because an output is declared in every path.
(In a real circuit, this is probably a designer error,
and a FIRRTL compiler is required to produce an error message.
The error message can be suppressed by a statement \lstinline!neg is invalid!.
Our formalization of the lowering transformation does not produce this error message
but leaves it to the LoFIRRTL compiler.)
The connection \lstinline~pos <= 0~ is overridden if \lstinline~condition~ is high;
this is translated to a multiplexer.

In the Coq implementation, the result of ``\texttt{find v m}'' (try to find the definition of \texttt{v} in map \texttt{m})
produces \texttt{None} if the component has not been declared;
\texttt{Some R\_default} if the component has been declared but not defined
(i.e.\@, not connected to);
and \texttt{Some (R\_fexpr e)} if the last connect assigned expression \texttt{e} to \texttt{v}.
At the end of a \lstinline!when! statement,
the maps of the two execution paths are joined into one map.
If the component has been declared in both paths%
---this means that the component has been declared before the \lstinline~when~ statement%
---but only defined in one, a \lstinline!validif! expression is produced.
If the component is defined in both paths
and the definitions differ, a \lstinline!mux! expression is produced.

The correctness proof compares the Coq implementation with (a formalization of) the output specified above.
It requires a mutual induction over statements and statement sequences
because the \lstinline~when~ statement contains two statement sequences,
but the proof is more long than complicated.
